%% abstract.tex

\begin{abstract}
Self-healing LLM agent operating systems increasingly adopt operating-system design principles---agents scheduled as processes, tools gated by permission checkers, and recovery orchestrators that mutate agent topology on failure. A classical OS invariant---\emph{an exception handler executes with the privileges of the faulting context, never higher}---has been enforced in kernels for decades, yet its transfer to the recovery path of LLM agent systems has not been examined. We show that this invariant is systematically violated: when a recovery orchestrator restructures the agent graph (e.g., replacing a failing agent's role), the new role is frequently selected by a diagnostic LLM whose input includes attacker-influenced content embedded in the failure payload, and the role-replacement sink performs no authorization check on the value it consumes. This is an instance of CWE-862 (Missing Authorization) at the recovery sink and the Confused Deputy problem at the orchestrator---two complementary framings of the same structural gap.

We approach this gap as a software architecture contribution, not a security measurement study. We formalize the \privinv{}---$P(\text{recovery action}) \leq P(\text{failed execution})$---as the agent-systems analog of the OS exception-handler privilege rule, and propose a design pattern that enforces it at a single centralized choke point: the \reauthgate{}. The pattern is realized by (i)~extending the recovery result type with a \code{requiresReauthorization} flag that side-effecting strategies set, (ii)~deferring the side effect as a callback stored in the recovery context rather than executing it inline, and (iii)~intercepting every flagged result at the orchestrator---the single component through which every recovery action must flow---and routing it through the existing \code{PermissionChecker} before the deferred callback executes. This localization turns the security property ``no recovery action exceeds the privilege of the failed execution'' from an emergent property of independently-authored strategies into a local, inspectable property of one orchestrator method.

We validate the pattern on \system{}, a multi-agent LLM operating system prototype (Java 21, 14 recovery strategy classes, four registered roles). Under the baseline configuration (no gate), a topology-mutation adversary that embeds role-escalation hints in a semantic core dump achieves a perfect 1.00 attack success rate across 100 trials, because the LLM-diagnosed role reaches the replacement sink without authorization. Under the defended configuration, the gate rejects every out-of-whitelist role \emph{by construction}---the invariant is a programmatic check, not a probabilistic defense---yielding 0.00 attack success, zero false alarms on benign recoveries, and sub-millisecond machinery overhead. Because the gate is LLM-independent, the guarantee is expected to hold across all models regardless of instruction-following strength.

The contribution is framed as a design pattern transfer: classical OS invariants (capability-based authority, exception-handler privilege containment) are ported to a domain---LLM agent recovery---where they have been conspicuously absent. The pattern is framework-agnostic in principle: it depends only on the existence of a recovery orchestrator and a permission checker, not on any specific agent architecture. The empirical validation establishes implementability and cost; a broader empirical evaluation (statistical generalization, cross-framework exploitability, adaptive adversaries) is beyond the scope of this architectural paper.
\end{abstract}
